\section{Conclusion}
\label{sec:conclusion}

We have shown that frontier LLMs systematically exhibit \emph{implicit avoidance} on gray zone topics --- legitimate questions involving social taboos, institutional critique, and epistemic uncertainty.
This avoidance manifests as hedging ($5.6\times$ higher than safe controls) and disclaimers ($3.5\times$ higher), not as explicit refusal (only 0.4\% of gray zone responses).
Avoidance patterns are highly correlated across four models from different organizations (mean Spearman $\rho = 0.78$), pointing to shared training corpora as the likely origin.

These findings have implications for AI safety research, deployment, and training.
Current overrefusal benchmarks miss implicit avoidance entirely because they measure binary refusal.
A new generation of benchmarks targeting hedging, qualification, and cautious persona behaviors is needed.
For deployment, the fact that models provide the least direct answers on precisely the topics where directness is most valuable --- contested scientific questions, institutional critique, sensitive social topics --- represents a systematic degradation of information quality.
For training, the cross-model consistency of avoidance patterns suggests that addressing this issue requires changes to training data composition, not just post-training safety tuning.

\para{Future work.}
Three directions are most promising.
First, activation-space analysis using open-weight models could identify an ``avoidance direction'' analogous to the refusal direction of \citet{zhao2025harmfulness}, enabling mechanistic understanding and targeted intervention.
Second, controlled finetuning experiments following \citet{betley2025emergent} could directly test whether narrow avoidance training produces broad avoidance generalization.
Third, human evaluation of implicit avoidance --- where annotators rate whether hedging is appropriate or excessive for each prompt --- would ground the phenomenon in user expectations rather than automated metrics.
